\documentclass[11pt]{article}

\usepackage[margin=1in]{geometry}
\usepackage{amsmath, amssymb}
\usepackage{graphicx}
\usepackage{enumitem}
\usepackage{hyperref}

\title{\textbf{Effect of Citronellal Odor Puffs on Cockroach Neurons}}
\author{Andres Rocha, Tyler Stoen, Jose Garcia}
\date{}

\begin{document}

\maketitle

%------------------------------------------------
\section*{Goal}

We study whether citronellal odor puff stimulation changes the firing rate of cockroach antennal lobe neurons. Specifically, we compare the response of three neurons and determine whether the stimulus effect is consistent across neurons. In this \href{https://github.com/tylerstoen/mp2/tree/main}{link} we have a repository containing code that produced the visuals along with a markdown file explaining it.

%------------------------------------------------
\section*{Exploratory Analysis and Hypotheses}

Raster plots show spike timing across trials. If stimulation affects
firing rate, spike density should change during the stimulus window
(5.99–6.49 seconds).

\begin{figure}[h]
\centering
\includegraphics[width=.32\textwidth]{WriteUp2ScreenShots/Screenshot 2026-03-04 at 09.23.12.png}
\includegraphics[width=.32\textwidth]{WriteUp2ScreenShots/Screenshot 2026-03-04 at 09.25.30.png}
\includegraphics[width=.32\textwidth]{WriteUp2ScreenShots/Screenshot 2026-03-04 at 09.27.22.png}
\caption{Raster plots for Neurons 1–3. The red band indicates the odor stimulus window.}
\end{figure}

\noindent Visual inspection suggests neuron-specific responses:

\begin{itemize}[noitemsep]
\item \textbf{Neuron 1:} Spike density increases during and immediately after the stimulus.
\item \textbf{Neuron 2:} Activity clusters around the stimulus period.
\item \textbf{Neuron 3:} Reduced spike activity during the stimulus window.
\end{itemize}

\noindent
\textbf{Hypotheses}

\begin{itemize}[noitemsep]
\item Neuron 1: stimulus increases firing rate
\item Neuron 2: stimulus increases firing rate
\item Neuron 3: stimulus suppresses firing rate
\end{itemize}

%------------------------------------------------
\section*{Model Specification}

We approximate the spike train using a non-homogeneous Poisson process
by binning time into small intervals:

\[
Y_{it} \sim \text{Poisson}(\lambda_i(t)\Delta t)
\]

\noindent Using a log link:

\[
\log(\lambda_i(t)) =
\beta_0 + \beta_1 \text{stim}(t) + f(t) + b_{\text{trial}}
\]

\noindent where

\begin{itemize}[noitemsep]
\item $\beta_1$ = stimulus effect (log scale)
\item $f(t)$ = smooth baseline firing pattern
\item $b_{\text{trial}} \sim N(0,\sigma^2)$ = trial random effect
\item $\log(\Delta t)$ = offset
\end{itemize}

\noindent This model allows firing intensity to vary smoothly in time while
estimating the multiplicative effect of the odor stimulus.

%------------------------------------------------
\section*{Results}

To interpret the stimulus effect we plot $\exp(\hat f(t))$, which represents the multiplicative change in firing intensity.

\begin{figure}[h]
\centering
\includegraphics[width=.32\textwidth]{WriteUp2ScreenShots/Screenshot 2026-03-04 at 09.40.17.png}
\includegraphics[width=.32\textwidth]{WriteUp2ScreenShots/Screenshot 2026-03-04 at 09.40.52.png}
\includegraphics[width=.32\textwidth]{WriteUp2ScreenShots/Screenshot 2026-03-04 at 09.41.27.png}
\caption{Estimated multiplicative stimulus effects for Neurons 1–3. The dashed line at 1 indicates baseline firing intensity. Values above 1 indicate increased firing, while values below 1 indicate suppression.}
\end{figure}

\begin{itemize}[noitemsep]
\item \textbf{Neuron 1:} The curve and confidence bands lie above 1 during
the stimulus window, indicating increased firing.

\item \textbf{Neuron 2:} The multiplicative effect rises above 1 during the
stimulus window, suggesting a moderate increase in firing.

\item \textbf{Neuron 3:} The curve falls below 1 during the stimulus,
indicating suppression of firing rate.
\end{itemize}

%------------------------------------------------
\section*{Conclusion}

\noindent \textbf{Do citronellal odor puffs affect firing rate?}

\noindent Yes. All neurons show stimulus effects that differ from baseline.

\smallskip

\noindent \textbf{Is the effect uniform across neurons?}

\smallskip

\noindent No. The response differs across neurons. Neurons 1 and 2 exhibit
increased firing during the stimulus, while Neuron 3 shows suppressed
activity.

\smallskip

\noindent These results suggest that citronellal odor stimulation produces
neuron-specific responses in the cockroach antennal lobe.

\end{document}
